\documentclass[12pt,a4paper]{article}
\usepackage[utf8]{vietnam}
\usepackage{amsmath,amssymb}
\usepackage{tabularx}
\usepackage{tikz}
\usepackage{xcolor}
\usepackage[top=0.6cm, bottom=0.6cm, left=0.6cm, right=0.6cm]{geometry}
\usepackage[dethi]{ex_test}

\SAOPTN{widthSA=2.6,heightSA=0.62}
\setlength{\parskip}{0pt}

% Dong ke cham mo chuan viet tay
\newcommand{\dongke}[1]{%
  \begingroup%
  \setlength{\parskip}{0pt}%
  \baselineskip=18.5pt%
  \vspace{1.0mm}%
  \foreach \i in {1,...,#1}{%
    \noindent{\color{gray!60}\dotfill}\par%
  }%
  \vspace{0.6mm}%
  \endgroup%
}

% Khu vuc bai lam chia 2 cot co vach ngan mo o giua
\newcommand{\dongkeHaiCot}[1]{%
  \par\vspace{0.5mm}%
  \noindent%
  \begin{minipage}[t]{0.485\linewidth}%
    \dongke{#1}%
  \end{minipage}%
  \hfill%
  {\color{gray!60}\vrule}%
  \hfill%
  \begin{minipage}[t]{0.485\linewidth}%
    \dongke{#1}%
  \end{minipage}%
}

% Dinh nghia o bai tap tu luan trong bang
\newcommand{\oBaiTap}[3][5]{%
  \vspace{1.0mm}%
  \noindent\textbf{#2} #3\par%
  \dongke{#1}%
}

\begin{document}
\pagestyle{empty}
% =============================================================================
% ========================= TOAN 9 TUAN 12 BUOI 02 - MA DE 101 =================
% ========================= TRANG 1/2 ==========================================
% =============================================================================
\noindent
\begin{tabularx}{\linewidth}{|X|c|}
\hline
\textbf{PHIẾU KIỂM TRA ĐỊNH KỲ TOÁN 9 - TUẦN 12 - BUỔI 02} & \textbf{Mã đề: 101 - Trang 1/2} \\
\hline
\end{tabularx}

\vspace{1.0mm}

\noindent
\begin{tabularx}{\linewidth}{|p{3.2cm}|X|}
\hline
\centering\textbf{ĐIỂM SỐ} & \centering\textbf{LỜI NHẬN XÉT CỦA THẦY CÔ} \tabularnewline
\hline
\rule{0pt}{1.8cm} & \tabularnewline
\hline
\end{tabularx}

\vspace{1.2mm}
\noindent\textbf{Câu 1 (0,25 điểm).} Cho $(O; 6\text{ cm})$, cung tròn có số đo $60^\circ$. Độ dài cung là: \shortans{2\pi\text{ cm}}\par\vspace{0.8mm}
\noindent\textbf{Câu 2 (0,25 điểm).} $\triangle ABC$ vuông tại $A$, $AB = 6\text{ cm}, AC = 8\text{ cm}$. Bán kính ngoại tiếp là: \shortans{5\text{ cm}}\par\vspace{0.8mm}
\noindent\textbf{Câu 3 (0,25 điểm).} Cho $(O; 5\text{ cm})$, dây $AB = 8\text{ cm}$. Khoảng cách từ $O$ đến $AB$ là: \shortans{3\text{ cm}}\par\vspace{0.8mm}
\noindent\textbf{Câu 4 (0,25 điểm).} Tính giá trị của biểu thức $\sqrt{(3 - \sqrt{10})^2} = $ \shortans{\sqrt{10} - 3}\par\vspace{0.8mm}
\noindent\textbf{Câu 5 (0,25 điểm).} Điều kiện xác định của biểu thức $\dfrac{1}{\sqrt{x - 2}}$ là: \shortans{x > 2}\par\vspace{0.8mm}
\noindent\textbf{Câu 6 (0,25 điểm).} Tính giá trị của biểu thức $\sqrt{75} - \sqrt{27} + \sqrt{12} = $ \shortans{4\sqrt{3}}\par\vspace{0.8mm}
\noindent\textbf{Câu 7 (0,25 điểm).} Hệ phương trình $\begin{cases} x + 2y = 8 \\ x - y = 2 \end{cases}$ có nghiệm $(x; y)$ là: \shortans{(4; 2)}\par\vspace{0.8mm}
\noindent\textbf{Câu 8 (0,25 điểm).} Tam giác vuông có hai cạnh góc vuông $6\text{ cm}, 8\text{ cm}$. Đường cao là: \shortans{4{,}8\text{ cm}}\par\vspace{0.8mm}
\noindent\textbf{Câu 9 (0,25 điểm).} Từ $A$ kẻ tiếp tuyến $AB$ với $(O; 5\text{ cm})$, biết $OA = 13\text{ cm}$. Độ dài $AB = $ \shortans{12\text{ cm}}\par\vspace{0.8mm}
\noindent\textbf{Câu 10 (0,25 điểm).} Bánh xe có bán kính $30\text{ cm}$, khi lăn $10$ vòng đi được: \shortans{6\pi\text{ m}}\par\vspace{0.8mm}
\vspace{0.6mm}
\noindent\textbf{Câu 11 (2,5 điểm).}
\vspace{0.6mm}

\noindent
\begin{tabularx}{\linewidth}{|X|X|}
\hline
\oBaiTap[19]{a) (1,25 điểm)}{Giải hệ phương trình:\par $\begin{cases} 2x - 3y = 1 \\ 3x + 2y = 8 \end{cases}$} & 
\oBaiTap[19]{b) (1,25 điểm)}{Cho biểu thức $P = \left(\dfrac{\sqrt{x}}{\sqrt{x} + 1} + \dfrac{2\sqrt{x}}{x - 1}\right) : \dfrac{\sqrt{x} + 1}{x - 1}$ với $x \ge 0, x \ne 1$.\par Rút gọn biểu thức $P$.} \\
\hline
\end{tabularx}

\newpage
% =============================================================================
% ========================= TOAN 9 TUAN 12 BUOI 02 - MA DE 101 =================
% ========================= TRANG 2/2 ==========================================
% =============================================================================
\noindent
\begin{tabularx}{\linewidth}{|X|c|}
\hline
\textbf{TOÁN 9 - TUẦN 12 - BUỔI 02} & \textbf{Mã đề: 101 - Trang 2/2} \\
\hline
\end{tabularx}

\vspace{1.2mm}

\noindent
\begin{tabularx}{\linewidth}{|X|}
\hline
\vspace{0.5mm}
\textbf{Câu 12 (2,5 điểm).} Theo kế hoạch, hai tổ công nhân phải may tổng cộng $800$ bộ quần áo. Thực tế, nhờ cải tiến kỹ thuật nên tổ I đã may vượt mức $15\%$, tổ II đã may vượt mức $20\%$ so với kế hoạch, do đó cả hai tổ đã may được tất cả $945$ bộ quần áo. Hỏi theo kế hoạch, mỗi tổ phải may bao nhiêu bộ quần áo?\par
\vspace{0.4mm}
\noindent\textit{Bài giải:}\dongkeHaiCot{15}
\\ \hline
\end{tabularx}

\vspace{2.5mm}

\noindent
\begin{tabularx}{\linewidth}{|X|}
\hline
\vspace{0.5mm}
\begin{minipage}[t]{0.54\linewidth}
\vspace{0pt}
\textbf{Câu 13 (2,5 điểm).} Cho đường tròn $(O; R)$ và điểm $A$ nằm bên ngoài đường tròn. Từ $A$ kẻ hai tiếp tuyến $AB, AC$ với đường tròn $(O)$ ($B, C$ là các tiếp điểm). Gọi $H$ là giao điểm của $OA$ và $BC$.\par\vspace{0.4mm}
\textbf{a) (1,25 điểm)} Chứng minh bốn điểm $A, B, O, C$ cùng thuộc một đường tròn.\par\vspace{0.4mm}
\textbf{b) (1,25 điểm)} Chứng minh $OA \perp BC$ tại $H$. Kẻ đường kính $BD$ của đường tròn $(O)$, chứng minh $CD \parallel OA$.
\end{minipage}\hfill
\begin{minipage}[t]{0.44\linewidth}
\vspace{0pt}\centering
\begin{tikzpicture}
  \draw[gray!60, dashed, rounded corners=3pt] (0, 0) rectangle (8.0, 3.3);
  \node[gray!60, below=1pt] at (4.0, 0) {\footnotesize\textit{(Khu vực vẽ hình)}};
\end{tikzpicture}
\end{minipage}\par
\vspace{0.4mm}
\noindent\textit{Bài giải:}\dongkeHaiCot{14}
\\ \hline
\end{tabularx}

\newpage
% =============================================================================
% ========================= TOAN 9 TUAN 12 BUOI 02 - MA DE 102 =================
% ========================= TRANG 1/2 ==========================================
% =============================================================================
\noindent
\begin{tabularx}{\linewidth}{|X|c|}
\hline
\textbf{PHIẾU KIỂM TRA ĐỊNH KỲ TOÁN 9 - TUẦN 12 - BUỔI 02} & \textbf{Mã đề: 102 - Trang 1/2} \\
\hline
\end{tabularx}

\vspace{1.0mm}

\noindent
\begin{tabularx}{\linewidth}{|p{3.2cm}|X|}
\hline
\centering\textbf{ĐIỂM SỐ} & \centering\textbf{LỜI NHẬN XÉT CỦA THẦY CÔ} \tabularnewline
\hline
\rule{0pt}{1.8cm} & \tabularnewline
\hline
\end{tabularx}

\vspace{1.2mm}
\noindent\textbf{Câu 1 (0,25 điểm).} Cho $(O; 9\text{ cm})$, cung tròn có số đo $40^\circ$. Độ dài cung là: \shortans{2\pi\text{ cm}}\par\vspace{0.8mm}
\noindent\textbf{Câu 2 (0,25 điểm).} $\triangle ABC$ vuông tại $A$, $AB = 9\text{ cm}, AC = 12\text{ cm}$. Bán kính ngoại tiếp là: \shortans{7{,}5\text{ cm}}\par\vspace{0.8mm}
\noindent\textbf{Câu 3 (0,25 điểm).} Cho $(O; 10\text{ cm})$, dây $AB = 16\text{ cm}$. Khoảng cách từ $O$ đến $AB$ là: \shortans{6\text{ cm}}\par\vspace{0.8mm}
\noindent\textbf{Câu 4 (0,25 điểm).} Tính giá trị của biểu thức $\sqrt{(4 - \sqrt{15})^2} = $ \shortans{4 - \sqrt{15}}\par\vspace{0.8mm}
\noindent\textbf{Câu 5 (0,25 điểm).} Điều kiện xác định của biểu thức $\dfrac{1}{\sqrt{x - 7}}$ là: \shortans{x > 7}\par\vspace{0.8mm}
\noindent\textbf{Câu 6 (0,25 điểm).} Tính giá trị của biểu thức $\sqrt{50} - \sqrt{18} + \sqrt{8} = $ \shortans{4\sqrt{2}}\par\vspace{0.8mm}
\noindent\textbf{Câu 7 (0,25 điểm).} Hệ phương trình $\begin{cases} 2x + y = 7 \\ x - y = 2 \end{cases}$ có nghiệm $(x; y)$ là: \shortans{(3; 1)}\par\vspace{0.8mm}
\noindent\textbf{Câu 8 (0,25 điểm).} Tam giác vuông có hai cạnh góc vuông $9\text{ cm}, 12\text{ cm}$. Đường cao là: \shortans{7{,}2\text{ cm}}\par\vspace{0.8mm}
\noindent\textbf{Câu 9 (0,25 điểm).} Từ $A$ kẻ tiếp tuyến $AB$ với $(O; 6\text{ cm})$, biết $OA = 10\text{ cm}$. Độ dài $AB = $ \shortans{8\text{ cm}}\par\vspace{0.8mm}
\noindent\textbf{Câu 10 (0,25 điểm).} Bánh xe có bán kính $35\text{ cm}$, khi lăn $10$ vòng đi được: \shortans{7\pi\text{ m}}\par\vspace{0.8mm}
\vspace{0.6mm}
\noindent\textbf{Câu 11 (2,5 điểm).}
\vspace{0.6mm}

\noindent
\begin{tabularx}{\linewidth}{|X|X|}
\hline
\oBaiTap[19]{a) (1,25 điểm)}{Giải hệ phương trình:\par $\begin{cases} 3x + 2y = 7 \\ 2x - 3y = -4 \end{cases}$} & 
\oBaiTap[19]{b) (1,25 điểm)}{Cho biểu thức $P = \left(\dfrac{\sqrt{x}}{\sqrt{x} + 2} + \dfrac{4\sqrt{x}}{x - 4}\right) : \dfrac{\sqrt{x} + 2}{x - 4}$ với $x \ge 0, x \ne 4$.\par Rút gọn biểu thức $P$.} \\
\hline
\end{tabularx}

\newpage
% =============================================================================
% ========================= TOAN 9 TUAN 12 BUOI 02 - MA DE 102 =================
% ========================= TRANG 2/2 ==========================================
% =============================================================================
\noindent
\begin{tabularx}{\linewidth}{|X|c|}
\hline
\textbf{TOÁN 9 - TUẦN 12 - BUỔI 02} & \textbf{Mã đề: 102 - Trang 2/2} \\
\hline
\end{tabularx}

\vspace{1.2mm}

\noindent
\begin{tabularx}{\linewidth}{|X|}
\hline
\vspace{0.5mm}
\textbf{Câu 12 (2,5 điểm).} Theo kế hoạch, hai tổ công nhân phải sản xuất tổng cộng $700$ chi tiết máy. Thực tế, tổ I làm vượt mức $15\%$, tổ II làm vượt mức $20\%$ so với kế hoạch, do đó cả hai tổ đã sản xuất được $820$ chi tiết máy. Hỏi theo kế hoạch, mỗi tổ phải sản xuất bao nhiêu chi tiết máy?\par
\vspace{0.4mm}
\noindent\textit{Bài giải:}\dongkeHaiCot{15}
\\ \hline
\end{tabularx}

\vspace{2.5mm}

\noindent
\begin{tabularx}{\linewidth}{|X|}
\hline
\vspace{0.5mm}
\begin{minipage}[t]{0.54\linewidth}
\vspace{0pt}
\textbf{Câu 13 (2,5 điểm).} Cho đường tròn $(O; R)$ và điểm $A$ nằm bên ngoài đường tròn. Từ $A$ kẻ hai tiếp tuyến $AB, AC$ với đường tròn $(O)$ ($B, C$ là các tiếp điểm). Gọi $H$ là giao điểm của $OA$ và $BC$.\par\vspace{0.4mm}
\textbf{a) (1,25 điểm)} Chứng minh bốn điểm $A, B, O, C$ cùng thuộc một đường tròn.\par\vspace{0.4mm}
\textbf{b) (1,25 điểm)} Chứng minh $OA \perp BC$ tại $H$. Kẻ đường kính $BD$ của đường tròn $(O)$, chứng minh $CD \parallel OA$.
\end{minipage}\hfill
\begin{minipage}[t]{0.44\linewidth}
\vspace{0pt}\centering
\begin{tikzpicture}
  \draw[gray!60, dashed, rounded corners=3pt] (0, 0) rectangle (8.0, 3.3);
  \node[gray!60, below=1pt] at (4.0, 0) {\footnotesize\textit{(Khu vực vẽ hình)}};
\end{tikzpicture}
\end{minipage}\par
\vspace{0.4mm}
\noindent\textit{Bài giải:}\dongkeHaiCot{14}
\\ \hline
\end{tabularx}

\newpage
% =============================================================================
% ========================= TOAN 9 TUAN 12 BUOI 02 - MA DE 103 =================
% ========================= TRANG 1/2 ==========================================
% =============================================================================
\noindent
\begin{tabularx}{\linewidth}{|X|c|}
\hline
\textbf{PHIẾU KIỂM TRA ĐỊNH KỲ TOÁN 9 - TUẦN 12 - BUỔI 02} & \textbf{Mã đề: 103 - Trang 1/2} \\
\hline
\end{tabularx}

\vspace{1.0mm}

\noindent
\begin{tabularx}{\linewidth}{|p{3.2cm}|X|}
\hline
\centering\textbf{ĐIỂM SỐ} & \centering\textbf{LỜI NHẬN XÉT CỦA THẦY CÔ} \tabularnewline
\hline
\rule{0pt}{1.8cm} & \tabularnewline
\hline
\end{tabularx}

\vspace{1.2mm}
\noindent\textbf{Câu 1 (0,25 điểm).} Cho $(O; 4\text{ cm})$, cung tròn có số đo $90^\circ$. Độ dài cung là: \shortans{2\pi\text{ cm}}\par\vspace{0.8mm}
\noindent\textbf{Câu 2 (0,25 điểm).} $\triangle ABC$ vuông tại $A$, $AB = 5\text{ cm}, AC = 12\text{ cm}$. Bán kính ngoại tiếp là: \shortans{6{,}5\text{ cm}}\par\vspace{0.8mm}
\noindent\textbf{Câu 3 (0,25 điểm).} Cho $(O; 13\text{ cm})$, dây $AB = 24\text{ cm}$. Khoảng cách từ $O$ đến $AB$ là: \shortans{5\text{ cm}}\par\vspace{0.8mm}
\noindent\textbf{Câu 4 (0,25 điểm).} Tính giá trị của biểu thức $\sqrt{(5 - \sqrt{27})^2} = $ \shortans{3\sqrt{3} - 5}\par\vspace{0.8mm}
\noindent\textbf{Câu 5 (0,25 điểm).} Điều kiện xác định của biểu thức $\dfrac{1}{\sqrt{x - 5}}$ là: \shortans{x > 5}\par\vspace{0.8mm}
\noindent\textbf{Câu 6 (0,25 điểm).} Tính giá trị của biểu thức $\sqrt{48} - \sqrt{12} + \sqrt{27} = $ \shortans{5\sqrt{3}}\par\vspace{0.8mm}
\noindent\textbf{Câu 7 (0,25 điểm).} Hệ phương trình $\begin{cases} 3x - y = 7 \\ x + y = 5 \end{cases}$ có nghiệm $(x; y)$ là: \shortans{(3; 2)}\par\vspace{0.8mm}
\noindent\textbf{Câu 8 (0,25 điểm).} Tam giác vuông có hai cạnh góc vuông $3\text{ cm}, 4\text{ cm}$. Đường cao là: \shortans{2{,}4\text{ cm}}\par\vspace{0.8mm}
\noindent\textbf{Câu 9 (0,25 điểm).} Từ $A$ kẻ tiếp tuyến $AB$ với $(O; 8\text{ cm})$, biết $OA = 17\text{ cm}$. Độ dài $AB = $ \shortans{15\text{ cm}}\par\vspace{0.8mm}
\noindent\textbf{Câu 10 (0,25 điểm).} Bánh xe có bán kính $25\text{ cm}$, khi lăn $10$ vòng đi được: \shortans{5\pi\text{ m}}\par\vspace{0.8mm}
\vspace{0.6mm}
\noindent\textbf{Câu 11 (2,5 điểm).}
\vspace{0.6mm}

\noindent
\begin{tabularx}{\linewidth}{|X|X|}
\hline
\oBaiTap[19]{a) (1,25 điểm)}{Giải hệ phương trình:\par $\begin{cases} 3x - 2y = 4 \\ 2x + 3y = 7 \end{cases}$} & 
\oBaiTap[19]{b) (1,25 điểm)}{Cho biểu thức $P = \left(\dfrac{\sqrt{x}}{\sqrt{x} + 3} + \dfrac{6\sqrt{x}}{x - 9}\right) : \dfrac{\sqrt{x} + 3}{x - 9}$ với $x \ge 0, x \ne 9$.\par Rút gọn biểu thức $P$.} \\
\hline
\end{tabularx}

\newpage
% =============================================================================
% ========================= TOAN 9 TUAN 12 BUOI 02 - MA DE 103 =================
% ========================= TRANG 2/2 ==========================================
% =============================================================================
\noindent
\begin{tabularx}{\linewidth}{|X|c|}
\hline
\textbf{TOÁN 9 - TUẦN 12 - BUỔI 02} & \textbf{Mã đề: 103 - Trang 2/2} \\
\hline
\end{tabularx}

\vspace{1.2mm}

\noindent
\begin{tabularx}{\linewidth}{|X|}
\hline
\vspace{0.5mm}
\textbf{Câu 12 (2,5 điểm).} Theo kế hoạch, hai tổ công nhân phải dệt tổng cộng $600$ chiếc áo len. Khi thực hiện, tổ I đã dệt vượt mức $10\%$, tổ II đã dệt vượt mức $15\%$ so với kế hoạch, do đó cả hai tổ đã dệt được tất cả $670$ chiếc áo len. Hỏi theo kế hoạch, mỗi tổ phải dệt bao nhiêu chiếc áo len?\par
\vspace{0.4mm}
\noindent\textit{Bài giải:}\dongkeHaiCot{15}
\\ \hline
\end{tabularx}

\vspace{2.5mm}

\noindent
\begin{tabularx}{\linewidth}{|X|}
\hline
\vspace{0.5mm}
\begin{minipage}[t]{0.54\linewidth}
\vspace{0pt}
\textbf{Câu 13 (2,5 điểm).} Cho đường tròn $(O; R)$ và điểm $A$ nằm bên ngoài đường tròn. Từ $A$ kẻ hai tiếp tuyến $AB, AC$ với đường tròn $(O)$ ($B, C$ là các tiếp điểm). Gọi $H$ là giao điểm của $OA$ và $BC$.\par\vspace{0.4mm}
\textbf{a) (1,25 điểm)} Chứng minh bốn điểm $A, B, O, C$ cùng thuộc một đường tròn.\par\vspace{0.4mm}
\textbf{b) (1,25 điểm)} Chứng minh $OA \perp BC$ tại $H$. Kẻ đường kính $BD$ của đường tròn $(O)$, chứng minh $CD \parallel OA$.
\end{minipage}\hfill
\begin{minipage}[t]{0.44\linewidth}
\vspace{0pt}\centering
\begin{tikzpicture}
  \draw[gray!60, dashed, rounded corners=3pt] (0, 0) rectangle (8.0, 3.3);
  \node[gray!60, below=1pt] at (4.0, 0) {\footnotesize\textit{(Khu vực vẽ hình)}};
\end{tikzpicture}
\end{minipage}\par
\vspace{0.4mm}
\noindent\textit{Bài giải:}\dongkeHaiCot{14}
\\ \hline
\end{tabularx}

\newpage
% =============================================================================
% ========================= TOAN 9 TUAN 12 BUOI 02 - MA DE 104 =================
% ========================= TRANG 1/2 ==========================================
% =============================================================================
\noindent
\begin{tabularx}{\linewidth}{|X|c|}
\hline
\textbf{PHIẾU KIỂM TRA ĐỊNH KỲ TOÁN 9 - TUẦN 12 - BUỔI 02} & \textbf{Mã đề: 104 - Trang 1/2} \\
\hline
\end{tabularx}

\vspace{1.0mm}

\noindent
\begin{tabularx}{\linewidth}{|p{3.2cm}|X|}
\hline
\centering\textbf{ĐIỂM SỐ} & \centering\textbf{LỜI NHẬN XÉT CỦA THẦY CÔ} \tabularnewline
\hline
\rule{0pt}{1.8cm} & \tabularnewline
\hline
\end{tabularx}

\vspace{1.2mm}
\noindent\textbf{Câu 1 (0,25 điểm).} Cho $(O; 10\text{ cm})$, cung tròn có số đo $36^\circ$. Độ dài cung là: \shortans{2\pi\text{ cm}}\par\vspace{0.8mm}
\noindent\textbf{Câu 2 (0,25 điểm).} $\triangle ABC$ vuông tại $A$, $AB = 8\text{ cm}, AC = 15\text{ cm}$. Bán kính ngoại tiếp là: \shortans{8{,}5\text{ cm}}\par\vspace{0.8mm}
\noindent\textbf{Câu 3 (0,25 điểm).} Cho $(O; 15\text{ cm})$, dây $AB = 24\text{ cm}$. Khoảng cách từ $O$ đến $AB$ là: \shortans{9\text{ cm}}\par\vspace{0.8mm}
\noindent\textbf{Câu 4 (0,25 điểm).} Tính giá trị của biểu thức $\sqrt{(2 - \sqrt{5})^2} = $ \shortans{\sqrt{5} - 2}\par\vspace{0.8mm}
\noindent\textbf{Câu 5 (0,25 điểm).} Điều kiện xác định của biểu thức $\dfrac{1}{\sqrt{x - 8}}$ là: \shortans{x > 8}\par\vspace{0.8mm}
\noindent\textbf{Câu 6 (0,25 điểm).} Tính giá trị của biểu thức $\sqrt{32} - \sqrt{8} + \sqrt{18} = $ \shortans{5\sqrt{2}}\par\vspace{0.8mm}
\noindent\textbf{Câu 7 (0,25 điểm).} Hệ phương trình $\begin{cases} x + y = 6 \\ 2x - y = 3 \end{cases}$ có nghiệm $(x; y)$ là: \shortans{(3; 3)}\par\vspace{0.8mm}
\noindent\textbf{Câu 8 (0,25 điểm).} Tam giác vuông có hai cạnh góc vuông $12\text{ cm}, 16\text{ cm}$. Đường cao là: \shortans{9{,}6\text{ cm}}\par\vspace{0.8mm}
\noindent\textbf{Câu 9 (0,25 điểm).} Từ $A$ kẻ tiếp tuyến $AB$ với $(O; 9\text{ cm})$, biết $OA = 15\text{ cm}$. Độ dài $AB = $ \shortans{12\text{ cm}}\par\vspace{0.8mm}
\noindent\textbf{Câu 10 (0,25 điểm).} Bánh xe có bán kính $40\text{ cm}$, khi lăn $10$ vòng đi được: \shortans{8\pi\text{ m}}\par\vspace{0.8mm}
\vspace{0.6mm}
\noindent\textbf{Câu 11 (2,5 điểm).}
\vspace{0.6mm}

\noindent
\begin{tabularx}{\linewidth}{|X|X|}
\hline
\oBaiTap[19]{a) (1,25 điểm)}{Giải hệ phương trình:\par $\begin{cases} 2x + 3y = 8 \\ 3x - 2y = -1 \end{cases}$} & 
\oBaiTap[19]{b) (1,25 điểm)}{Cho biểu thức $P = \left(\dfrac{\sqrt{x}}{\sqrt{x} + 4} + \dfrac{8\sqrt{x}}{x - 16}\right) : \dfrac{\sqrt{x} + 4}{x - 16}$ với $x \ge 0, x \ne 16$.\par Rút gọn biểu thức $P$.} \\
\hline
\end{tabularx}

\newpage
% =============================================================================
% ========================= TOAN 9 TUAN 12 BUOI 02 - MA DE 104 =================
% ========================= TRANG 2/2 ==========================================
% =============================================================================
\noindent
\begin{tabularx}{\linewidth}{|X|c|}
\hline
\textbf{TOÁN 9 - TUẦN 12 - BUỔI 02} & \textbf{Mã đề: 104 - Trang 2/2} \\
\hline
\end{tabularx}

\vspace{1.2mm}

\noindent
\begin{tabularx}{\linewidth}{|X|}
\hline
\vspace{0.5mm}
\textbf{Câu 12 (2,5 điểm).} Theo kế hoạch, hai đội công nhân phải trồng tổng cộng $500$ cây xanh. Trong đợt thi đua, đội I đã trồng vượt mức $15\%$, đội II đã trồng vượt mức $10\%$ so với kế hoạch, do đó cả hai đội đã trồng được tất cả $565$ cây xanh. Hỏi theo kế hoạch, mỗi đội phải trồng bao nhiêu cây xanh?\par
\vspace{0.4mm}
\noindent\textit{Bài giải:}\dongkeHaiCot{15}
\\ \hline
\end{tabularx}

\vspace{2.5mm}

\noindent
\begin{tabularx}{\linewidth}{|X|}
\hline
\vspace{0.5mm}
\begin{minipage}[t]{0.54\linewidth}
\vspace{0pt}
\textbf{Câu 13 (2,5 điểm).} Cho đường tròn $(O; R)$ và điểm $A$ nằm bên ngoài đường tròn. Từ $A$ kẻ hai tiếp tuyến $AB, AC$ với đường tròn $(O)$ ($B, C$ là các tiếp điểm). Gọi $H$ là giao điểm của $OA$ và $BC$.\par\vspace{0.4mm}
\textbf{a) (1,25 điểm)} Chứng minh bốn điểm $A, B, O, C$ cùng thuộc một đường tròn.\par\vspace{0.4mm}
\textbf{b) (1,25 điểm)} Chứng minh $OA \perp BC$ tại $H$. Kẻ đường kính $BD$ của đường tròn $(O)$, chứng minh $CD \parallel OA$.
\end{minipage}\hfill
\begin{minipage}[t]{0.44\linewidth}
\vspace{0pt}\centering
\begin{tikzpicture}
  \draw[gray!60, dashed, rounded corners=3pt] (0, 0) rectangle (8.0, 3.3);
  \node[gray!60, below=1pt] at (4.0, 0) {\footnotesize\textit{(Khu vực vẽ hình)}};
\end{tikzpicture}
\end{minipage}\par
\vspace{0.4mm}
\noindent\textit{Bài giải:}\dongkeHaiCot{14}
\\ \hline
\end{tabularx}
\end{document}
